\section{Le débat fondationnel : trois positions classiques}
\label{sec:debat}

Depuis l'article fondateur d'Aharonov et Bohm
\cite{AharonovBohm1959}, l'électrodynamique quantique pose une
question fondationnelle qui résiste à toute réduction à
l'expérience : que doit-on considérer comme l'objet physiquement
réel parmi le 4-potentiel $A_\mu$, le tenseur de champ
$F_{\mu\nu}$, et la phase d'holonomie
$\oint_\gamma A \cdot d\ell$ ? Trois positions se partagent la
littérature contemporaine, sans qu'aucune expérience classique
n'ait pu les départager. Cette section les rappelle dans leur
forme la plus tranchée, dans le seul but de fixer ce que la
Théorie de la Persistance (PT) prétend trancher. La revue n'a pas
vocation à être exhaustive --- pour les développements
philosophiques on renvoie aux monographies de Healey
\cite{Healey2007} et Maudlin~\cite{Maudlin2011}.

\subsection{Position A : la réalité du potentiel vecteur}
\label{ssec:position_A}

La position A, héritée du programme de Bohm--Hiley
\cite{BohmHiley1993}, soutient que le potentiel $A_\mu$ est
physiquement réel et que le tenseur $F_{\mu\nu} = \partial_\mu
A_\nu - \partial_\nu A_\mu$ en est une grandeur dérivée. L'argument
empirique est l'effet Aharonov--Bohm \cite{AharonovBohm1959} :
le déphasage interférométrique
$\Delta\varphi = (e/\hbar) \oint A \cdot d\ell$ enregistré par une
particule qui contourne un solénoïde idéal est mesurable même là
où $E = B = 0$ sur tout son trajet. Aucun champ ne traverse le
trajet ; aucun objet local ne peut donc rendre compte de l'effet.
Le potentiel, lui, est non nul (à jauge près) le long du trajet,
et le programme bohmien y voit la preuve qu'il faut le considérer
comme une grandeur d'ontologie supérieure.

La difficulté connue de la position A est la dépendance de jauge.
Le potentiel $A_\mu$ n'est défini qu'à une transformation
$A_\mu \mapsto A_\mu + \partial_\mu \Lambda$ près, et toute
quantité physique mesurable doit être invariante sous cette
transformation. Or la phase Aharonov--Bohm n'est invariante de
jauge \emph{que dans son intégrale fermée} ; ses valeurs locales
ne le sont pas. La position A doit donc soit accepter une réalité
physique qui n'est pas elle-même invariante de jauge --- ce que
beaucoup considèrent comme une concession ontologique majeure ---
soit traiter $A_\mu$ comme un objet hybride dont seule la classe
d'équivalence sous transformations de jauge serait réelle.

\subsection{Position B : la non-localité intrinsèque du champ}
\label{ssec:position_B}

La position B accepte $F_{\mu\nu}$ comme l'objet réel mais en
fait un objet intrinsèquement non-local. Dans la lecture forte
proposée par Mead \cite{Mead2000} et dans le sillage des analyses
de Bell~\cite{Bell1964} et des tests d'Aspect \cite{Aspect1982} et
Hensen et al.~\cite{Hensen2015}, le champ
électromagnétique n'est pas une distribution sur l'espace-temps
classique : il porte une structure de corrélation entre points
spatialement séparés qui ne s'explique pas par la variation
locale de ses composantes. La phase Aharonov--Bohm devient alors
le signal d'une telle non-localité, et les corrélations EPR celui
de son extension à la matière.

L'objection structurale à la position B est qu'elle exige de
modifier l'ontologie classique du champ pour la rendre
intrinsèquement non-locale, sans préciser quel mécanisme
mathématique implémente cette modification. Les approches
formelles disponibles (fibrés de connexions, théories de jauge,
opérateurs de Wilson) sont toutes interprétables en termes
strictement locaux moyennant un choix adapté de variables. La
non-localité de la position B reste donc, dans la majorité des
formulations, une propriété de l'\emph{interprétation} du champ
plutôt que de sa structure mathématique.

\subsection{Position C : la primauté de l'holonomie}
\label{ssec:position_C}

La position C, due à Wu et Yang \cite{WuYang1975} et reprise par
Healey \cite{Healey2007} et Belot \cite{Belot1998}, soutient
qu'aucun des deux objets locaux $A_\mu, F_{\mu\nu}$ n'est
fondamental : $A_\mu$ parce qu'il dépend de la jauge,
$F_{\mu\nu}$ parce qu'il ne contient pas l'information de la
phase Aharonov--Bohm dans les régions où il s'annule. Le seul
objet à la fois jauge-invariant et sensible à l'expérience
Aharonov--Bohm est la boucle de Wilson
$W_\gamma = \exp\bigl[i (e/\hbar) \oint_\gamma A \cdot d\ell\bigr]$,
fonctionnelle non-locale des chemins fermés~$\gamma$. C'est elle,
et non $A_\mu$ ou $F_{\mu\nu}$, qui mérite le statut d'objet
primitif.

La position C est conceptuellement satisfaisante --- elle
dissout l'asymétrie entre la jauge-dépendance de $A_\mu$ et
l'incomplétude de $F_{\mu\nu}$ --- mais elle laisse ouverte la
question de l'\emph{origine} des boucles de Wilson : pourquoi
ces objets non-locaux particuliers, et non d'autres? La théorie
de jauge classique les introduit comme une construction
mathématique post hoc, à partir d'une connexion qui reste
elle-même primaire dans la formulation. La position C reste donc,
en un sens, une description et non une explication.

\subsection{Le constat opérationnel}
\label{ssec:constat}

Les trois positions A, B, C font les mêmes prédictions
expérimentales dans le cadre de l'électrodynamique quantique
standard. La phase Aharonov--Bohm est mesurable dans les trois
cadres ; les corrélations EPR/Bell sont reproduites dans les
trois ; les sections efficaces et les durées de vie sont
identiques. Aucune mesure n'a pu, à ce jour, départager les trois
positions au niveau de la prédiction quantitative. Le débat est
donc, à l'heure actuelle, un débat d'interprétation.

La Théorie de la Persistance modifie cet état de fait sur deux
points. D'une part, elle promeut la position~C en \emph{théorème}
plutôt qu'en interprétation : l'angle d'holonomie modulaire
$\theta_p$ est l'unique invariant primitif satisfaisant quatre
conditions structurales que ni $A_\mu$ ni $F_{\mu\nu}$ ne peuvent
remplir simultanément (Section~\ref{sec:holonomie_primitive}).
D'autre part, elle produit une prédiction expérimentale --- la
prédiction QH1a de la Section~\ref{sec:prediction_qh1a} --- dont
le verdict empirique distingue effectivement la PT (et avec elle
la position~C) d'une théorie standard sur variété courbe. Le
débat fondationnel n'est plus seulement un débat
d'interprétation : il est, sous PT, un débat tranchable.
